% --- Archon physics formalization source begin ---
% archon:physics
% archon:covers PhyXMiniProblems/problem_phyx_mini_0330.lean
% archon:source-report reports/phyx_mini/problem_phyx_mini_0330.source.json

\chapter{Physics problem phyx\_mini\_0330}
\label{ch:PhyXMiniProblems_problem_phyx_mini_0330}

\paragraph{Problem source.}
\#\# Physical scenario

In figure, a point source \$S\$ of sound waves lies near a reflecting wall \$AB\$. A sound detector \$D\$ intercepts sound ray \$R\_1\$ traveling directly from \$S\$. It also intercepts sound ray \$R\_2\$ that reflects from the wall such that the angle of incidence \$\textbackslash{}theta\_i\$ is equal to the angle of reflection \$\textbackslash{}theta\_r\$. Assume that the reflection of sound by the wall causes a phase shift of \$0.500\textbackslash{}lambda\$.

\#\# Question

If the distances are \$d\_1 = 2.50\$ m, \$d\_2 = 20.0\$ m, and \$d\_3 = 12.5\$ m, what is the second lowest frequency at which \$R\_1\$ and \$R\_2\$ are in phase at \$D\$?

\#\# Answer choices

- A: A: 98 Hz

- B: B: 102 Hz

- C: C: 109 Hz

- D: D: 118 Hz

\#\# Figure caption (auxiliary; use the image as primary evidence)

The image is a diagram illustrating a reflection scenario involving lines, angles, and distances. 

- **Objects and Labels**:
  - A vertical wall on the left, labeled as section "AB."
  - Two points are marked as "S" and "D." 
  - Point "S" is close to the wall, and point "D" is further away in the lower right section.
  - The wall is detailed with a textured pattern.

- **Lines**:
  - A horizontal dashed line extends from the wall to point "D," passing through point "S." This line is divided into segments labeled \textbackslash{}(d\_1\textbackslash{}) and \textbackslash{}(d\_2\textbackslash{}).
  - A vertical dashed line drops from the end of the dashed horizontal line to connect with point "D," labeled as \textbackslash{}(d\_3\textbackslash{}).
  - Two red lines connect point "S" to point "D." These lines represent paths labeled \textbackslash{}(R\_1\textbackslash{}) and \textbackslash{}(R\_2\textbackslash{}).

- **Angles**:
  - Angles \textbackslash{}(\textbackslash{}theta\_i\textbackslash{}) and \textbackslash{}(\textbackslash{}theta\_r\textbackslash{}) are formed at the point "S" near the wall, illustrating the angle of incidence and reflection relative to the line perpendicular to the wall.
  
- **Relationships**:
  - The lines and angles suggest a reflection concept, with light or sound originating from point "S" and reflecting off the wall toward point "D."
  - The red lines, angles, and distances together represent a geometrical analysis typically used in optics or acoustics.

This combination of lines, angles, and labels likely represents a study or analysis of reflection principles using a geometric approach.

\#\# Dataset metadata

Resonance and Harmonics; Spatial Relation Reasoning, Multi-Formula Reasoning

\paragraph{Recorded answer/context.}
D: D: 118 Hz

\paragraph{Figure/image path.}
/root/proposal\_for\_physic/science-mango/phyx\_mini\_run/phyx\_data/test\_image/330.png

\paragraph{Formalization target.}
create a compiling Lean file with sorry bodies at `PhyXMiniProblems/problem\_phyx\_mini\_0330.lean`. The Lean declarations must preserve the physical quantities, dimensions or dimensional roles, figure labels, governing-law hypotheses, and final relation expressed by this problem.
Use Mathlib/Physlib names found through LeanExplore where available. If a physics API is missing, introduce faithful local abstractions rather than scalar placeholder aliases.

\begin{theorem}[Physics formalization target]
\label{thm:physics:phyx_mini_0330:target}
\lean{PhyXMiniProblems.ProblemPhyXMini0330.problem_phyx_mini_0330}
\uses{def:physics:phyx-mini-0330:phyxminiproblems-problemphyxmini0330-acousticfrequency, def:physics:phyx-mini-0330:phyxminiproblems-problemphyxmini0330-frequencyinhertz, def:physics:phyx-mini-0330:phyxminiproblems-problemphyxmini0330-speedinmeterspersecond, def:physics:phyx-mini-0330:phyxminiproblems-problemphyxmini0330-wallreflectedsoundsetup, def:physics:phyx-mini-0330:phyxminiproblems-problemphyxmini0330-matchessuppliedfigure, def:physics:phyx-mini-0330:phyxminiproblems-problemphyxmini0330-matchesproblemreadouts, def:physics:phyx-mini-0330:phyxminiproblems-problemphyxmini0330-hasphysicalacousticparameters, def:physics:phyx-mini-0330:phyxminiproblems-problemphyxmini0330-satisfiesspecularreflection, def:physics:phyx-mini-0330:phyxminiproblems-problemphyxmini0330-satisfieswallreflectionpathgeometry, def:physics:phyx-mini-0330:phyxminiproblems-problemphyxmini0330-satisfiespathdifferencelaw, def:physics:phyx-mini-0330:phyxminiproblems-problemphyxmini0330-satisfiessoundspeedfrequencywavelengthlaw, def:physics:phyx-mini-0330:phyxminiproblems-problemphyxmini0330-satisfiesreflectedsoundinterferencelaw, def:physics:phyx-mini-0330:phyxminiproblems-problemphyxmini0330-issecondlowestinphasefrequency}
The assigned autoformalize agent should translate this physics problem into Lean declarations in the covered file, with theorem and lemma proof bodies written as `by sorry`.
\end{theorem}
\begin{proof}
This is an autoformalization task, not a proof task. Produce statements that can later be proved by the physics prover without weakening the physical model.
\end{proof}

% --- Archon declaration topology begin ---
\section{Declaration topology}
\begin{definition}[Acoustic Length]
  \label{def:physics:phyx-mini-0330:phyxminiproblems-problemphyxmini0330-acousticlength}
  \lean{PhyXMiniProblems.ProblemPhyXMini0330.AcousticLength}
  A nonnegative physical length, independent of the unit used to read it
\end{definition}
\begin{proof}
  This is the declared model definition of Acoustic Length.
\end{proof}
\begin{definition}[Acoustic Frequency]
  \label{def:physics:phyx-mini-0330:phyxminiproblems-problemphyxmini0330-acousticfrequency}
  \lean{PhyXMiniProblems.ProblemPhyXMini0330.AcousticFrequency}
  A nonnegative cyclic frequency, carrying the inverse-time dimension
\end{definition}
\begin{proof}
  This is the declared model definition of Acoustic Frequency.
\end{proof}
\begin{definition}[Acoustic Speed]
  \label{def:physics:phyx-mini-0330:phyxminiproblems-problemphyxmini0330-acousticspeed}
  \lean{PhyXMiniProblems.ProblemPhyXMini0330.AcousticSpeed}
  A nonnegative physical propagation speed
\end{definition}
\begin{proof}
  This is the declared model definition of Acoustic Speed.
\end{proof}
\begin{definition}[length Readout]
  \label{def:physics:phyx-mini-0330:phyxminiproblems-problemphyxmini0330-lengthreadout}
  \lean{PhyXMiniProblems.ProblemPhyXMini0330.lengthReadout}
  \uses{def:physics:phyx-mini-0330:phyxminiproblems-problemphyxmini0330-acousticlength}
  Read a physical length in the selected length unit
\end{definition}
\begin{proof}
  This is the declared model definition of length Readout.
\end{proof}
\begin{definition}[frequency Readout]
  \label{def:physics:phyx-mini-0330:phyxminiproblems-problemphyxmini0330-frequencyreadout}
  \lean{PhyXMiniProblems.ProblemPhyXMini0330.frequencyReadout}
  \uses{def:physics:phyx-mini-0330:phyxminiproblems-problemphyxmini0330-acousticfrequency}
  Read a physical cyclic frequency in inverse units of the selected time unit
\end{definition}
\begin{proof}
  This is the declared model definition of frequency Readout.
\end{proof}
\begin{definition}[speed Readout]
  \label{def:physics:phyx-mini-0330:phyxminiproblems-problemphyxmini0330-speedreadout}
  \lean{PhyXMiniProblems.ProblemPhyXMini0330.speedReadout}
  \uses{def:physics:phyx-mini-0330:phyxminiproblems-problemphyxmini0330-acousticspeed}
  Read a speed in the selected length unit per selected time unit
\end{definition}
\begin{proof}
  This is the declared model definition of speed Readout.
\end{proof}
\begin{definition}[length In Meters]
  \label{def:physics:phyx-mini-0330:phyxminiproblems-problemphyxmini0330-lengthinmeters}
  \lean{PhyXMiniProblems.ProblemPhyXMini0330.lengthInMeters}
  \uses{def:physics:phyx-mini-0330:phyxminiproblems-problemphyxmini0330-acousticlength, def:physics:phyx-mini-0330:phyxminiproblems-problemphyxmini0330-lengthreadout}
  Metre readout used for the three distances and the ray path lengths
\end{definition}
\begin{proof}
  This is the declared model definition of length In Meters.
\end{proof}
\begin{definition}[frequency In Hertz]
  \label{def:physics:phyx-mini-0330:phyxminiproblems-problemphyxmini0330-frequencyinhertz}
  \lean{PhyXMiniProblems.ProblemPhyXMini0330.frequencyInHertz}
  \uses{def:physics:phyx-mini-0330:phyxminiproblems-problemphyxmini0330-acousticfrequency, def:physics:phyx-mini-0330:phyxminiproblems-problemphyxmini0330-frequencyreadout}
  Hertz readout, i.e. cycles per SI second
\end{definition}
\begin{proof}
  This is the declared model definition of frequency In Hertz.
\end{proof}
\begin{definition}[speed In Meters Per Second]
  \label{def:physics:phyx-mini-0330:phyxminiproblems-problemphyxmini0330-speedinmeterspersecond}
  \lean{PhyXMiniProblems.ProblemPhyXMini0330.speedInMetersPerSecond}
  \uses{def:physics:phyx-mini-0330:phyxminiproblems-problemphyxmini0330-acousticspeed, def:physics:phyx-mini-0330:phyxminiproblems-problemphyxmini0330-speedreadout}
  Metre-per-second readout of the sound speed
\end{definition}
\begin{proof}
  This is the declared model definition of speed In Meters Per Second.
\end{proof}
\begin{definition}[Figure Point]
  \label{def:physics:phyx-mini-0330:phyxminiproblems-problemphyxmini0330-figurepoint}
  \lean{PhyXMiniProblems.ProblemPhyXMini0330.FigurePoint}
  The four literal point labels printed in the primary figure
\end{definition}
\begin{proof}
  This is the declared model definition of Figure Point.
\end{proof}
\begin{definition}[Wall Label]
  \label{def:physics:phyx-mini-0330:phyxminiproblems-problemphyxmini0330-walllabel}
  \lean{PhyXMiniProblems.ProblemPhyXMini0330.WallLabel}
  The label of the reflecting wall segment
\end{definition}
\begin{proof}
  This is the declared model definition of Wall Label.
\end{proof}
\begin{definition}[Wall Orientation]
  \label{def:physics:phyx-mini-0330:phyxminiproblems-problemphyxmini0330-wallorientation}
  \lean{PhyXMiniProblems.ProblemPhyXMini0330.WallOrientation}
  The orientation of wall AB in the supplied figure
\end{definition}
\begin{proof}
  This is the declared model definition of Wall Orientation.
\end{proof}
\begin{definition}[Sound Ray]
  \label{def:physics:phyx-mini-0330:phyxminiproblems-problemphyxmini0330-soundray}
  \lean{PhyXMiniProblems.ProblemPhyXMini0330.SoundRay}
  The two red sound-ray labels printed in the figure
\end{definition}
\begin{proof}
  This is the declared model definition of Sound Ray.
\end{proof}
\begin{definition}[Ray Route Kind]
  \label{def:physics:phyx-mini-0330:phyxminiproblems-problemphyxmini0330-rayroutekind}
  \lean{PhyXMiniProblems.ProblemPhyXMini0330.RayRouteKind}
  The geometric behavior of a depicted sound ray
\end{definition}
\begin{proof}
  This is the declared model definition of Ray Route Kind.
\end{proof}
\begin{definition}[Detector Placement]
  \label{def:physics:phyx-mini-0330:phyxminiproblems-problemphyxmini0330-detectorplacement}
  \lean{PhyXMiniProblems.ProblemPhyXMini0330.DetectorPlacement}
  The relative placement of D visible in the figure
\end{definition}
\begin{proof}
  This is the declared model definition of Detector Placement.
\end{proof}
\begin{definition}[Acoustic Medium]
  \label{def:physics:phyx-mini-0330:phyxminiproblems-problemphyxmini0330-acousticmedium}
  \lean{PhyXMiniProblems.ProblemPhyXMini0330.AcousticMedium}
  The propagation medium implicit in the sound-reflection experiment
\end{definition}
\begin{proof}
  This is the declared model definition of Acoustic Medium.
\end{proof}
\begin{definition}[Depicted Sound Ray]
  \label{def:physics:phyx-mini-0330:phyxminiproblems-problemphyxmini0330-depictedsoundray}
  \lean{PhyXMiniProblems.ProblemPhyXMini0330.DepictedSoundRay}
  \uses{def:physics:phyx-mini-0330:phyxminiproblems-problemphyxmini0330-figurepoint, def:physics:phyx-mini-0330:phyxminiproblems-problemphyxmini0330-rayroutekind}
  A red ray drawn from the labelled source to the labelled detector
\end{definition}
\begin{proof}
  This is the declared model definition of Depicted Sound Ray.
\end{proof}
\begin{definition}[Sound Reflection Figure]
  \label{def:physics:phyx-mini-0330:phyxminiproblems-problemphyxmini0330-soundreflectionfigure}
  \lean{PhyXMiniProblems.ProblemPhyXMini0330.SoundReflectionFigure}
  \uses{def:physics:phyx-mini-0330:phyxminiproblems-problemphyxmini0330-figurepoint, def:physics:phyx-mini-0330:phyxminiproblems-problemphyxmini0330-walllabel, def:physics:phyx-mini-0330:phyxminiproblems-problemphyxmini0330-wallorientation, def:physics:phyx-mini-0330:phyxminiproblems-problemphyxmini0330-soundray, def:physics:phyx-mini-0330:phyxminiproblems-problemphyxmini0330-detectorplacement, def:physics:phyx-mini-0330:phyxminiproblems-problemphyxmini0330-depictedsoundray}
  Qualitative information and dimensionless angles read from the primary image. The incidence and reflection angles belong to the unique wall encounter of R2; their equality is imposed later as the law of specular reflection
\end{definition}
\begin{proof}
  This is the declared model definition of Sound Reflection Figure.
\end{proof}
\begin{definition}[Wall Reflected Sound Setup]
  \label{def:physics:phyx-mini-0330:phyxminiproblems-problemphyxmini0330-wallreflectedsoundsetup}
  \lean{PhyXMiniProblems.ProblemPhyXMini0330.WallReflectedSoundSetup}
  \uses{def:physics:phyx-mini-0330:phyxminiproblems-problemphyxmini0330-acousticlength, def:physics:phyx-mini-0330:phyxminiproblems-problemphyxmini0330-acousticfrequency, def:physics:phyx-mini-0330:phyxminiproblems-problemphyxmini0330-acousticspeed, def:physics:phyx-mini-0330:phyxminiproblems-problemphyxmini0330-soundray, def:physics:phyx-mini-0330:phyxminiproblems-problemphyxmini0330-acousticmedium, def:physics:phyx-mini-0330:phyxminiproblems-problemphyxmini0330-soundreflectionfigure}
  The independent physical quantities of the experiment. d1 is the perpendicular horizontal distance from wall AB to source S; d2 is the horizontal separation from S to the vertical line through D; and d3 is the vertical separation down to D. rayPathLength, pathDifference, wavelengthAtFrequency, and inPhaseAtDetector are independent physical observables constrained only by the governing laws below
\end{definition}
\begin{proof}
  This is the declared model definition of Wall Reflected Sound Setup.
\end{proof}
\begin{definition}[Matches Supplied Figure]
  \label{def:physics:phyx-mini-0330:phyxminiproblems-problemphyxmini0330-matchessuppliedfigure}
  \lean{PhyXMiniProblems.ProblemPhyXMini0330.MatchesSuppliedFigure}
  \uses{def:physics:phyx-mini-0330:phyxminiproblems-problemphyxmini0330-wallreflectedsoundsetup}
  Labels, routes, and relative placement transcribed from the primary bitmap. Both rays start at S and reach D; R1 is direct and R2 reflects once from wall AB. No path length, frequency, or interference conclusion is contained here
\end{definition}
\begin{proof}
  This is the declared model definition of Matches Supplied Figure.
\end{proof}
\begin{definition}[Matches Problem Readouts]
  \label{def:physics:phyx-mini-0330:phyxminiproblems-problemphyxmini0330-matchesproblemreadouts}
  \lean{PhyXMiniProblems.ProblemPhyXMini0330.MatchesProblemReadouts}
  \uses{def:physics:phyx-mini-0330:phyxminiproblems-problemphyxmini0330-lengthinmeters, def:physics:phyx-mini-0330:phyxminiproblems-problemphyxmini0330-wallreflectedsoundsetup}
  Numerical data printed in the problem. The terminating decimals are represented exactly as rational real readouts. The stated 0.500 lambda reflection shift is the dimensionless fraction 1/2; it is not a wavelength or frequency answer
\end{definition}
\begin{proof}
  This is the declared model definition of Matches Problem Readouts.
\end{proof}
\begin{definition}[Uses Standard Air Sound Speed]
  \label{def:physics:phyx-mini-0330:phyxminiproblems-problemphyxmini0330-usesstandardairsoundspeed}
  \lean{PhyXMiniProblems.ProblemPhyXMini0330.UsesStandardAirSoundSpeed}
  \uses{def:physics:phyx-mini-0330:phyxminiproblems-problemphyxmini0330-speedinmeterspersecond, def:physics:phyx-mini-0330:phyxminiproblems-problemphyxmini0330-wallreflectedsoundsetup}
  The problem does not print a propagation speed. This separate calibration records the conventional room-temperature value 343 m/s used to evaluate the whole-hertz answer choices
\end{definition}
\begin{proof}
  This is the declared model definition of Uses Standard Air Sound Speed.
\end{proof}
\begin{definition}[Has Physical Acoustic Parameters]
  \label{def:physics:phyx-mini-0330:phyxminiproblems-problemphyxmini0330-hasphysicalacousticparameters}
  \lean{PhyXMiniProblems.ProblemPhyXMini0330.HasPhysicalAcousticParameters}
  \uses{def:physics:phyx-mini-0330:phyxminiproblems-problemphyxmini0330-lengthreadout, def:physics:phyx-mini-0330:phyxminiproblems-problemphyxmini0330-speedreadout, def:physics:phyx-mini-0330:phyxminiproblems-problemphyxmini0330-frequencyinhertz, def:physics:phyx-mini-0330:phyxminiproblems-problemphyxmini0330-wallreflectedsoundsetup}
  Positivity and nondegeneracy conditions for the physical setup
\end{definition}
\begin{proof}
  This is the declared model definition of Has Physical Acoustic Parameters.
\end{proof}
\begin{definition}[Satisfies Specular Reflection]
  \label{def:physics:phyx-mini-0330:phyxminiproblems-problemphyxmini0330-satisfiesspecularreflection}
  \lean{PhyXMiniProblems.ProblemPhyXMini0330.SatisfiesSpecularReflection}
  \uses{def:physics:phyx-mini-0330:phyxminiproblems-problemphyxmini0330-wallreflectedsoundsetup}
  The reflected ray obeys the equal-angle law at wall AB
\end{definition}
\begin{proof}
  This is the declared model definition of Satisfies Specular Reflection.
\end{proof}
\begin{definition}[Satisfies Wall Reflection Path Geometry]
  \label{def:physics:phyx-mini-0330:phyxminiproblems-problemphyxmini0330-satisfieswallreflectionpathgeometry}
  \lean{PhyXMiniProblems.ProblemPhyXMini0330.SatisfiesWallReflectionPathGeometry}
  \uses{def:physics:phyx-mini-0330:phyxminiproblems-problemphyxmini0330-lengthreadout, def:physics:phyx-mini-0330:phyxminiproblems-problemphyxmini0330-wallreflectedsoundsetup}
  The direct length follows from the d2 by d3 right triangle. For the reflected length, mirroring S across the vertical wall changes the horizontal leg to 2*d1 + d2; the unfolded ray is then another right-triangle hypotenuse. These are general figure-geometry equations in any length unit
\end{definition}
\begin{proof}
  This is the declared model definition of Satisfies Wall Reflection Path Geometry.
\end{proof}
\begin{definition}[Satisfies Path Difference Law]
  \label{def:physics:phyx-mini-0330:phyxminiproblems-problemphyxmini0330-satisfiespathdifferencelaw}
  \lean{PhyXMiniProblems.ProblemPhyXMini0330.SatisfiesPathDifferenceLaw}
  \uses{def:physics:phyx-mini-0330:phyxminiproblems-problemphyxmini0330-lengthreadout, def:physics:phyx-mini-0330:phyxminiproblems-problemphyxmini0330-wallreflectedsoundsetup}
  The path difference is the magnitude of the two ray-length difference
\end{definition}
\begin{proof}
  This is the declared model definition of Satisfies Path Difference Law.
\end{proof}
\begin{definition}[Satisfies Sound Speed Frequency Wavelength Law]
  \label{def:physics:phyx-mini-0330:phyxminiproblems-problemphyxmini0330-satisfiessoundspeedfrequencywavelengthlaw}
  \lean{PhyXMiniProblems.ProblemPhyXMini0330.SatisfiesSoundSpeedFrequencyWavelengthLaw}
  \uses{def:physics:phyx-mini-0330:phyxminiproblems-problemphyxmini0330-lengthreadout, def:physics:phyx-mini-0330:phyxminiproblems-problemphyxmini0330-frequencyreadout, def:physics:phyx-mini-0330:phyxminiproblems-problemphyxmini0330-speedreadout, def:physics:phyx-mini-0330:phyxminiproblems-problemphyxmini0330-frequencyinhertz, def:physics:phyx-mini-0330:phyxminiproblems-problemphyxmini0330-wallreflectedsoundsetup}
  The nondispersive sound-wave relation v = lambda*f, stated in compatible arbitrary length and time units. It characterizes wavelength at every positive frequency without assigning the requested frequency
\end{definition}
\begin{proof}
  This is the declared model definition of Satisfies Sound Speed Frequency Wavelength Law.
\end{proof}
\begin{definition}[Satisfies Reflected Sound Interference Law]
  \label{def:physics:phyx-mini-0330:phyxminiproblems-problemphyxmini0330-satisfiesreflectedsoundinterferencelaw}
  \lean{PhyXMiniProblems.ProblemPhyXMini0330.SatisfiesReflectedSoundInterferenceLaw}
  \uses{def:physics:phyx-mini-0330:phyxminiproblems-problemphyxmini0330-lengthreadout, def:physics:phyx-mini-0330:phyxminiproblems-problemphyxmini0330-frequencyinhertz, def:physics:phyx-mini-0330:phyxminiproblems-problemphyxmini0330-wallreflectedsoundsetup}
  For coherent rays from one point source, arrival is in phase exactly when the propagation path difference plus the reflection phase displacement is an integral number of wavelengths. The positive natural order ranges over the entire constructive-interference spectrum, so this law does not select the first, second, or any numerical frequency
\end{definition}
\begin{proof}
  This is the declared model definition of Satisfies Reflected Sound Interference Law.
\end{proof}
\begin{definition}[Is Positive In Phase Frequency]
  \label{def:physics:phyx-mini-0330:phyxminiproblems-problemphyxmini0330-ispositiveinphasefrequency}
  \lean{PhyXMiniProblems.ProblemPhyXMini0330.IsPositiveInPhaseFrequency}
  \uses{def:physics:phyx-mini-0330:phyxminiproblems-problemphyxmini0330-acousticfrequency, def:physics:phyx-mini-0330:phyxminiproblems-problemphyxmini0330-frequencyinhertz, def:physics:phyx-mini-0330:phyxminiproblems-problemphyxmini0330-wallreflectedsoundsetup}
  A positive physical frequency at which the two received rays are in phase
\end{definition}
\begin{proof}
  This is the declared model definition of Is Positive In Phase Frequency.
\end{proof}
\begin{definition}[Is Lowest Positive In Phase Frequency]
  \label{def:physics:phyx-mini-0330:phyxminiproblems-problemphyxmini0330-islowestpositiveinphasefrequency}
  \lean{PhyXMiniProblems.ProblemPhyXMini0330.IsLowestPositiveInPhaseFrequency}
  \uses{def:physics:phyx-mini-0330:phyxminiproblems-problemphyxmini0330-acousticfrequency, def:physics:phyx-mini-0330:phyxminiproblems-problemphyxmini0330-frequencyinhertz, def:physics:phyx-mini-0330:phyxminiproblems-problemphyxmini0330-wallreflectedsoundsetup, def:physics:phyx-mini-0330:phyxminiproblems-problemphyxmini0330-ispositiveinphasefrequency}
  A least positive in-phase frequency, expressed through its hertz order
\end{definition}
\begin{proof}
  This is the declared model definition of Is Lowest Positive In Phase Frequency.
\end{proof}
\begin{definition}[Is Second Lowest In Phase Frequency]
  \label{def:physics:phyx-mini-0330:phyxminiproblems-problemphyxmini0330-issecondlowestinphasefrequency}
  \lean{PhyXMiniProblems.ProblemPhyXMini0330.IsSecondLowestInPhaseFrequency}
  \uses{def:physics:phyx-mini-0330:phyxminiproblems-problemphyxmini0330-acousticfrequency, def:physics:phyx-mini-0330:phyxminiproblems-problemphyxmini0330-frequencyinhertz, def:physics:phyx-mini-0330:phyxminiproblems-problemphyxmini0330-wallreflectedsoundsetup, def:physics:phyx-mini-0330:phyxminiproblems-problemphyxmini0330-ispositiveinphasefrequency, def:physics:phyx-mini-0330:phyxminiproblems-problemphyxmini0330-islowestpositiveinphasefrequency}
  frequency is second-lowest when a lowest positive in-phase frequency lies strictly below it and it is no greater than any positive in-phase frequency strictly above that lowest one. This generic definition contains no numerical frequency, interference order, or answer label
\end{definition}
\begin{proof}
  This is the declared model definition of Is Second Lowest In Phase Frequency.
\end{proof}
\begin{definition}[Answer Choice]
  \label{def:physics:phyx-mini-0330:phyxminiproblems-problemphyxmini0330-answerchoice}
  \lean{PhyXMiniProblems.ProblemPhyXMini0330.AnswerChoice}
  The four answer labels printed with the problem
\end{definition}
\begin{proof}
  This is the declared model definition of Answer Choice.
\end{proof}
\begin{definition}[hertz]
  \label{def:physics:phyx-mini-0330:phyxminiproblems-problemphyxmini0330-answerchoice-hertz}
  \lean{PhyXMiniProblems.ProblemPhyXMini0330.AnswerChoice.hertz}
  \uses{def:physics:phyx-mini-0330:phyxminiproblems-problemphyxmini0330-answerchoice}
  Whole-hertz value displayed beside each answer label
\end{definition}
\begin{proof}
  This is the declared model definition of hertz.
\end{proof}
\begin{definition}[recorded Dataset Answer]
  \label{def:physics:phyx-mini-0330:phyxminiproblems-problemphyxmini0330-recordeddatasetanswer}
  \lean{PhyXMiniProblems.ProblemPhyXMini0330.recordedDatasetAnswer}
  \uses{def:physics:phyx-mini-0330:phyxminiproblems-problemphyxmini0330-answerchoice}
  The answer label recorded in the supplied dataset metadata
\end{definition}
\begin{proof}
  This is the declared model definition of recorded Dataset Answer.
\end{proof}
\begin{definition}[Rounds To Whole Hertz]
  \label{def:physics:phyx-mini-0330:phyxminiproblems-problemphyxmini0330-roundstowholehertz}
  \lean{PhyXMiniProblems.ProblemPhyXMini0330.RoundsToWholeHertz}
  \uses{def:physics:phyx-mini-0330:phyxminiproblems-problemphyxmini0330-acousticfrequency, def:physics:phyx-mini-0330:phyxminiproblems-problemphyxmini0330-frequencyinhertz}
  A physical frequency rounds to a displayed whole-hertz value
\end{definition}
\begin{proof}
  This is the declared model definition of Rounds To Whole Hertz.
\end{proof}
\begin{definition}[Agrees With Displayed Frequency Choice]
  \label{def:physics:phyx-mini-0330:phyxminiproblems-problemphyxmini0330-agreeswithdisplayedfrequencychoice}
  \lean{PhyXMiniProblems.ProblemPhyXMini0330.AgreesWithDisplayedFrequencyChoice}
  \uses{def:physics:phyx-mini-0330:phyxminiproblems-problemphyxmini0330-acousticfrequency, def:physics:phyx-mini-0330:phyxminiproblems-problemphyxmini0330-answerchoice, def:physics:phyx-mini-0330:phyxminiproblems-problemphyxmini0330-roundstowholehertz}
  A physical frequency agrees, after rounding, with a displayed choice
\end{definition}
\begin{proof}
  This is the declared model definition of Agrees With Displayed Frequency Choice.
\end{proof}
\begin{lemma}[path Difference In Meters exact]
  \label{lem:physics:phyx-mini-0330:phyxminiproblems-problemphyxmini0330-pathdifferenceinmeters-exact}
  \lean{PhyXMiniProblems.ProblemPhyXMini0330.pathDifferenceInMeters_exact}
  \uses{def:physics:phyx-mini-0330:phyxminiproblems-problemphyxmini0330-lengthinmeters, def:physics:phyx-mini-0330:phyxminiproblems-problemphyxmini0330-wallreflectedsoundsetup, def:physics:phyx-mini-0330:phyxminiproblems-problemphyxmini0330-matchesproblemreadouts, def:physics:phyx-mini-0330:phyxminiproblems-problemphyxmini0330-hasphysicalacousticparameters, def:physics:phyx-mini-0330:phyxminiproblems-problemphyxmini0330-satisfieswallreflectionpathgeometry, def:physics:phyx-mini-0330:phyxminiproblems-problemphyxmini0330-satisfiespathdifferencelaw}
  The image-source geometry and the supplied distance readouts give the exact SI path difference sqrt (25\textasciicircum{}2 + (25/2)\textasciicircum{}2) - sqrt (20\textasciicircum{}2 + (25/2)\textasciicircum{}2) metres. This helper is a geometric consequence, not a frequency conclusion
\end{lemma}
\begin{proof}
  The conclusion follows from the governing relations and figure readouts named in the statement of path Difference In Meters exact.
\end{proof}
\begin{lemma}[second Lowest In Phase Frequency formula]
  \label{lem:physics:phyx-mini-0330:phyxminiproblems-problemphyxmini0330-secondlowestinphasefrequency-formula}
  \lean{PhyXMiniProblems.ProblemPhyXMini0330.secondLowestInPhaseFrequency_formula}
  \uses{def:physics:phyx-mini-0330:phyxminiproblems-problemphyxmini0330-acousticfrequency, def:physics:phyx-mini-0330:phyxminiproblems-problemphyxmini0330-lengthinmeters, def:physics:phyx-mini-0330:phyxminiproblems-problemphyxmini0330-frequencyinhertz, def:physics:phyx-mini-0330:phyxminiproblems-problemphyxmini0330-speedinmeterspersecond, def:physics:phyx-mini-0330:phyxminiproblems-problemphyxmini0330-wallreflectedsoundsetup, def:physics:phyx-mini-0330:phyxminiproblems-problemphyxmini0330-matchesproblemreadouts, def:physics:phyx-mini-0330:phyxminiproblems-problemphyxmini0330-hasphysicalacousticparameters, def:physics:phyx-mini-0330:phyxminiproblems-problemphyxmini0330-satisfieswallreflectionpathgeometry, def:physics:phyx-mini-0330:phyxminiproblems-problemphyxmini0330-satisfiespathdifferencelaw, def:physics:phyx-mini-0330:phyxminiproblems-problemphyxmini0330-satisfiessoundspeedfrequencywavelengthlaw, def:physics:phyx-mini-0330:phyxminiproblems-problemphyxmini0330-satisfiesreflectedsoundinterferencelaw, def:physics:phyx-mini-0330:phyxminiproblems-problemphyxmini0330-issecondlowestinphasefrequency}
  With a half-wavelength reflection phase shift, consecutive in-phase orders have Delta L = (order - 1/2) lambda. Thus the second positive order has f = 3 v / (2 Delta L). The existential frequency is still a dimensionful physical quantity; this lemma does not round it or mention an answer choice
\end{lemma}
\begin{proof}
  The conclusion follows from the governing relations and figure readouts named in the statement of second Lowest In Phase Frequency formula.
\end{proof}
\begin{theorem}[problem phyx mini 0330 at Standard Air Speed]
  \label{thm:physics:phyx-mini-0330:phyxminiproblems-problemphyxmini0330-problem-phyx-mini-0330-atstandardairspeed}
  \lean{PhyXMiniProblems.ProblemPhyXMini0330.problem_phyx_mini_0330_atStandardAirSpeed}
  \uses{def:physics:phyx-mini-0330:phyxminiproblems-problemphyxmini0330-acousticfrequency, def:physics:phyx-mini-0330:phyxminiproblems-problemphyxmini0330-frequencyinhertz, def:physics:phyx-mini-0330:phyxminiproblems-problemphyxmini0330-wallreflectedsoundsetup, def:physics:phyx-mini-0330:phyxminiproblems-problemphyxmini0330-matchessuppliedfigure, def:physics:phyx-mini-0330:phyxminiproblems-problemphyxmini0330-matchesproblemreadouts, def:physics:phyx-mini-0330:phyxminiproblems-problemphyxmini0330-usesstandardairsoundspeed, def:physics:phyx-mini-0330:phyxminiproblems-problemphyxmini0330-hasphysicalacousticparameters, def:physics:phyx-mini-0330:phyxminiproblems-problemphyxmini0330-satisfiesspecularreflection, def:physics:phyx-mini-0330:phyxminiproblems-problemphyxmini0330-satisfieswallreflectionpathgeometry, def:physics:phyx-mini-0330:phyxminiproblems-problemphyxmini0330-satisfiespathdifferencelaw, def:physics:phyx-mini-0330:phyxminiproblems-problemphyxmini0330-satisfiessoundspeedfrequencywavelengthlaw, def:physics:phyx-mini-0330:phyxminiproblems-problemphyxmini0330-satisfiesreflectedsoundinterferencelaw, def:physics:phyx-mini-0330:phyxminiproblems-problemphyxmini0330-issecondlowestinphasefrequency, def:physics:phyx-mini-0330:phyxminiproblems-problemphyxmini0330-recordeddatasetanswer, def:physics:phyx-mini-0330:phyxminiproblems-problemphyxmini0330-roundstowholehertz, def:physics:phyx-mini-0330:phyxminiproblems-problemphyxmini0330-agreeswithdisplayedfrequencychoice}
  If one separately adopts the conventional room-air calibration 343 m/s, the symbolic result is approximately 117.845 Hz. It therefore rounds to 118 Hz and agrees with displayed and recorded choice D. This is deliberately a conditional corollary rather than part of the source-faithful target
\end{theorem}
\begin{proof}
  The conclusion follows from the governing relations and figure readouts named in the statement of problem phyx mini 0330 at Standard Air Speed.
\end{proof}
% --- Archon declaration topology end ---
% --- Archon physics formalization source end ---
